% Regression: regress_x_pbc — renders B^i = X A^i X^{-1} with open virtual stubs and on-wire X capsules.
% Formula: Gauge equation B^i = X A^i X^{-1}: both sides are D x D matrices with
\documentclass[varwidth=16cm, border=6pt]{standalone}
\usepackage{amsmath, amssymb}
\usepackage{tenkz}
\begin{document}

% Gauge equation B^i = X A^i X^{-1}: both sides are D x D matrices with
% one physical index -- open virtual stubs on both sides, X and X^{-1}
% as on-wire matrix capsules.
G:
\[
  \begin{tenkz}[physical=up]
    \tn[up=$i$]{B}
  \end{tenkz}
  \;=\;
  \begin{tenkz}[physical=up]
    \tnX{X} & \tn[up=$i$]{A} & \tnX{X^{-1}}
  \end{tenkz}
\]

% Tr(A^i B^j): periodic returns the row's east end to its west end --
% the virtual index is summed, i and j stay open.
T:
\begin{tenkz}[periodic, physical=up]
  \tn[up=$i$]{A} & \tn[up=$j$]{B}
\end{tenkz}

\medskip
% Same closure with the return wire on the label side: dot labels south,
% legs north -- exercises the label-band addition below the row.
T2:
\begin{tenkz}[periodic]
  \tn{A} & \tn{B} & \tn{C}
\end{tenkz}

\medskip
% Tr over both virtual rows of a two-row word: the above-racetrack must
% clear the labelled up legs of the ket row -- the tip-label scan adds
% the label band on top of the leg length.
T3:
\begin{tenkz}[rows={ket:nopair, bra}, periodic]
  \tn[up=$i$]{A} & \tn[up=$j$]{A} \\
  \tn*{A} & \tn*{A}
\end{tenkz}

\end{document}
